%======================================================================
\subsection{BW-ROU next: connectors, fittings, drawings, the harness as a vault item (plan, 30 Oct 2026)}\label{sec:rou-next-plan}
%======================================================================

On 30 Oct you asked for all four follow-ups to the routing track (\S\ref{sec:session-30oct}). Two of them, the
Assembly BOM \emph{Length} column and the configurable bend rules, need no design and are being built now
(\S\ref{sec:rou-built-30oct}). The other three change the data model, so they come here first with their
questions. They build in this order, because each one needs the one before it:

\begin{center}
\small
\begin{tabular}{@{}lll@{}}
\toprule
Step & What & Needs\\
\midrule
1 & Connectors (pins) and fittings (ports): \eref{ROU-05} & ROU-01..04\\
2 & The harness as a vault item, the purchasing BOM by length: \eref{ROU-06} & step 1, PDM-04 vault\\
3 & Drawings: flattened harness, tube with bend table: \eref{ROU-07} & steps 1--2, DRW-01\\
\bottomrule
\end{tabular}
\end{center}

%----------------------------------------------------------------------
\subsubsection*{Step 1: connectors and fittings}

\paragraph{The idea.} A connector or fitting is an ordinary part (a vault item, often COTS) with
\textbf{named attachment frames} in it: \code{Pin1}, \code{Pin2}, \dots{} on a connector and \code{Port1},
\dots{} on a fitting. Each frame is an LCS whose origin is where the conductor or tube enters and whose $+Z$
points out along the entry direction. The frames carry a few properties:
\begin{center}
\small
\begin{tabularx}{\linewidth}{@{}lX@{}}
\toprule
Pin & name (``1'', ``A3'', ``VBAT''), accepted gauge range (e.g.\ 28--22\,AWG), terminal P/N (crimp contact),
      seal P/N, default signal\\
Port & port type (push-to-connect, barb, compression, NPT/BSPP thread), tube or hose size it takes, minimum
      straight length before the first bend\\
\bottomrule
\end{tabularx}
\end{center}
Routing then gets three behaviours: (1) a route path's first or last point picked on a pin or port
\emph{attaches} to that frame (ROU-01 already attaches to LCS origins); (2) the path \emph{leaves along the
frame's $+Z$} for the minimum straight length, so a wire does not bend at the crimp and a tube enters a fitting
straight; (3) a wire on that path fills its \code{From}/\code{FromPin} (\code{To}/\code{ToPin}) from the
connector's label and the pin's name, and the from-to list stops being typed by hand. Mating pairs (a plug and
its receptacle) record each other's P/N, so the from-to list can show a connection through a mated pair.

\paragraph{Where the geometry comes from.} Most robotics and drone connectors (JST-XH/PH/GH/SH, Molex
Micro-Fit, XT30/XT60, Dupont 2.54\,mm headers) already have STEP models in KiCad's 3D library. The KiCad
bridge (\eref{ECD-05}) already brings KiCad parts into the Standard Libraries vault, and the footprint's pads
give the pin positions and names. A connector made this way is a KiCad-library item plus the pin frames
generated from its pads (question C1). Fittings (push-to-connect, barbs, NPT elbows) have no such library.
They are made as parametric parts from a size table, as the tube sizes are in ROU-03.

\paragraph{Terminals, seals and sleeving.} A crimped wire end needs a terminal (and in sealed connectors a
seal). When a wire attaches to a pin with a terminal P/N, the terminal appears in the BOM as a counted line
under the harness (question C3). Heat-shrink and braided sleeving are cut-to-length items on a bundle and go
to the BOM by length, like wire.

%----------------------------------------------------------------------
\subsubsection*{Step 2: the harness as a vault item}

\paragraph{The problem.} You chose in R5 (\S\ref{sec:rou}) that a harness is its own vault item with a harness
P/N, because it is bought or built as one thing, with one P/N and one drawing. But it is
routed \emph{in context}: its paths attach to the motor, the board and the frame in the parent assembly. In
FreeCAD the cleanest way to keep both is the SolidWorks pattern of a \emph{virtual} component that is saved out
later:
\begin{enumerate}[nosep]
  \item a \textbf{Harness} container (an \code{App::Part} with a \code{HarnessType} marker) holds the paths,
        wires, bundles, sleeving and the connector instances, inside the assembly file, while you route;
  \item \textbf{Make vault item} saves it out as its own file under a new P/N (document type \emph{Cable
        assembly} / \emph{Tube assembly}). The parent links to it, and the path points keep their attachments
        as cross-document links (\code{App::PropertyXLink}, which FreeCAD resolves when both files are
        open, as in-context editing \eref{ASM-07} does);
  \item opened on its own, the harness keeps the last-solved point positions, so its drawing and BOM do not
        need the parent.
\end{enumerate}

\paragraph{The purchasing BOM.} Wire, tube, hose and sleeving become \textbf{stock items}: one vault item per
kind, e.g.\ ``Wire, 22\,AWG, red, PVC hook-up'', with \code{QuantityUnit}\,=\,m (question V2). Every routed
item references its stock item. The harness's own BOM then lists its connectors and terminals by count and
its stock items by length. The flattened vault BOM \emph{sums} the lengths per stock item across all harnesses
and adds a scrap allowance, the rounding to a purchase unit (a 30\,m spool) being optional:
\[
  L_{\text{buy}}=\Bigl\lceil \frac{(1+s)\sum_k n_k L_k}{L_{\text{unit}}}\Bigr\rceil L_{\text{unit}},
\]
with $n_k$ the number of harnesses of kind $k$ in the product and $L_k$ the length of that stock item in one
of them.

%----------------------------------------------------------------------
\subsubsection*{Step 3: drawings}

\paragraph{Flattened harness drawing.} The harness is a tree: connectors are leaves, breakouts are inner
nodes and bundle segments are edges with lengths. The drawing lays that tree out flat in the sheet's plane
(figure~\ref{fig:rou-flat}). It keeps the topology and labels every segment with its true length, but it is
\emph{not to scale}, as most harness drawings are (question D1). The layout is a straightforward tree drawing:
the longest path (trunk) runs horizontally, branches leave at $\pm45^\circ$ alternating sides, and a
force-directed pass separates labels that overlap. The sheet carries a connector table (connector, P/N,
mating P/N, a pin-out table per connector), the from-to wire list, the stock-item list with lengths, and
notes (lengths $\pm$ tolerance, labels, sleeving).

\begin{figure}[htbp]
\centering
\begin{tikzpicture}[x=1mm,y=1mm,font=\sffamily\scriptsize]
  \draw[line width=1.6pt, black!70] (0,20) -- (40,20) -- (85,20);
  \draw[line width=1.2pt, black!70] (40,20) -- (55,35);
  \draw[line width=1.2pt, black!70] (40,20) -- (55,5);
  \foreach \p/\n in {{(0,20)}/J1, {(85,20)}/J2, {(55,35)}/M1, {(55,5)}/S1}{
    \fill[bwblue] \p ++(-2.5,-2.5) rectangle ++(5,5);
    \node[above=2.2mm] at \p {\n};}
  \fill[black] (40,20) circle (0.9);
  \node[below=1.2mm] at (20,20) {180 $\pm$5};
  \node[below=1.2mm] at (63,20) {220 $\pm$5};
  \node[anchor=east] at (47,30) {90};
  \node[anchor=east] at (47,10) {120};
  \node[anchor=west, text=black!55] at (90,30) {J1 board, J2 battery};
  \node[anchor=west, text=black!55] at (90,25) {M1 motor, S1 sensor};
  \node[anchor=west, text=black!55] at (90,18) {segment lengths true, drawing not to scale};
  \node[anchor=west, text=black!55] at (90,13) {breakout ($\bullet$) dimensioned from J1};
\end{tikzpicture}
\caption{A flattened harness: the routed 3D tree laid out flat, trunk horizontal, branches at $45^\circ$,
every segment labelled with its routed length. Connector, pin-out and wire tables sit beside it.}
\label{fig:rou-flat}
\end{figure}

\paragraph{Tube drawing with bend table.} A tube bender works from the \textbf{YBC} (also called LRA) table:
feed length $Y$, rotation $B$ and bend angle $C$ per bend. ROU-03 already computes it. With the intersection
points $\mathbf P_0\dots\mathbf P_{n+1}$ of the polyline, the unit legs
$\mathbf d_i=(\mathbf P_{i+1}-\mathbf P_i)/\lVert\cdot\rVert$, the bend radius $R$ and $t_i=R\tan(C_i/2)$:
\[
  C_i=\arccos(\mathbf d_{i-1}\cdot\mathbf d_i),\qquad
  Y_i=\lVert\mathbf P_{i}-\mathbf P_{i-1}\rVert-t_{i-1}-t_i,\qquad
  B_i=\angle\bigl(\mathbf d_{i-2}\times\mathbf d_{i-1},\ \mathbf d_{i-1}\times\mathbf d_i\bigr),
\]
where $B_i$ is the signed dihedral angle between successive bend planes about $\mathbf d_{i-1}$ (with $t_0=0$
at the start). The developed (cut) length is $\sum Y_i+\sum R\,C_i$ plus the last straight. The drawing has an
isometric and a top view of the tube, the YBC table, the XYZ intersection-point table (what CNC benders also
accept), and the tube size, material, CLR and cut length in the title block notes. Springback and elongation
compensation depend on the bender and the material, so they are left out for now (question D3).

%----------------------------------------------------------------------
\subsubsection*{New epics}
\begin{itemize}[nosep]
  \item \eref{ROU-05} Connectors and fittings: pin and port frames, attach and leave straight, auto from-to,
        terminals in the BOM, starter library.
  \item \eref{ROU-06} Harness as a vault item: harness container, Make vault item, stock items, purchasing BOM
        by length with scrap and rounding.
  \item \eref{ROU-07} Routing drawings: flattened harness drawing with tables, tube drawing with YBC and XYZ
        tables.
\end{itemize}

\paragraph{Questions for you.}
\begin{enumerate}[label=C\arabic*.]
  \item \textbf{Connector geometry.} KiCad 3D models with pins generated from the footprint pads (through
        the ECD-05 bridge), parametric outlines we draw ourselves, or manufacturer STEP files you download?
        \emph{Recommendation:} KiCad models and pads; parametric only where KiCad has none.
  \item \textbf{Starter library.} Which families first? Proposed: JST-XH, JST-PH, JST-GH, JST-SH, Dupont
        2.54\,mm, Molex Micro-Fit 3.0, XT30, XT60, XT90, MR30, 3.5\,mm bullet (motor leads), Deutsch DT.
  \item \textbf{Terminals in the BOM} automatically from the pin's terminal P/N (one per attached wire end)?
        \emph{Recommendation:} yes.
\end{enumerate}
\begin{enumerate}[label=F\arabic*.]
  \item \textbf{Fitting families first.} Push-to-connect (4/6/8\,mm pneumatic), hose barbs, NPT/BSPP threaded
        elbows and tees, compression? \emph{Recommendation:} push-to-connect, barbs and NPT/BSPP; compression
        later.
  \item \textbf{Straight length at a pin or port.} A minimum straight length before the first bend, set on the
        pin or port (defaults: wire $2\times$OD past the terminal, tube $2\times$OD or the bender's clamp
        length), configurable like the bend rules? \emph{Recommendation:} yes.
\end{enumerate}
\begin{enumerate}[label=V\arabic*.]
  \item \textbf{Harness model.} How to reach R5's own vault item: route in context as a harness container,
        then \emph{Make vault item} saves it out (above), or each harness is its own file from the start and is
        edited in context?
        \emph{Recommendation:} the container first, saved out later.
  \item \textbf{Stock items.} Wire, tube, hose and sleeving as vault items with their own P/N per gauge,
        colour and insulation (the purchasing BOM sums by P/N), or grouped by properties without a P/N?
        \emph{Recommendation:} vault items with P/N, document type \emph{Stock material}, shared with
        \eref{PDM-11}'s stock forms.
  \item \textbf{Scrap and rounding.} A scrap allowance (default 10\,\%) and optional rounding to the purchase
        unit (spool length) per stock item? \emph{Recommendation:} yes, both per stock item, scrap default
        10\,\%.
\end{enumerate}
\begin{enumerate}[label=D\arabic*.]
  \item \textbf{Harness drawing scale.} Not to scale with every segment dimensioned (the usual practice), or
        a 1:1 formboard print (tiled PDF) as well? \emph{Recommendation:} not to scale now; the 1:1 formboard
        is a later option.
  \item \textbf{Harness drawing content.} Connector table with pin-outs, from-to list, stock-item lengths,
        labels and notes as above? Anything your builds need beyond that (wire labels/markers, test table)?
  \item \textbf{Tube drawing.} YBC and XYZ tables; springback compensation later?
        \emph{Recommendation:} yes, springback later.
\end{enumerate}
